% !TEX root = ../main.tex

\chapter{Background}
\label{ch:background}
This chapter builds up the pieces the thesis depends on, ending with PDDLStream and how 3D space is represented for planning.

\section{AI Planning Fundamentals}
\Ac{ai} planning produces a sequence of actions---a plan---that gets from a given initial state to a specified goal \cite{ghallab2004automated}.

\subsection{State-Space Models in Classical Planning}
Classical planning describes problems in a \textit{factored} representation: states are specified through the atoms that hold and actions through their effects on those atoms, rather than by enumerating every configuration of the world explicitly, which is infeasible for all but the smallest problems \cite{geffner2013concise, ghallab2004automated}.

A foundational language for this is \textbf{\acs{strips}} \cite{fikes1971strips}. A \textbf{state} is the set of ground atoms (propositions) true in a configuration of the world---for example, in a blocksworld with blocks A, B and a table, $\{\texttt{(on A B)}, \texttt{(ontable B)}, \texttt{(handempty)}\}$. An \textbf{action} $a$ is given by three sets of atoms: preconditions $\mathrm{pre}(a)$, an add-list $\mathrm{add}(a)$, and a delete-list $\mathrm{del}(a)$. Action $a$ is \textit{applicable} in a state $s$ iff $\mathrm{pre}(a) \subseteq s$, and applying it produces the successor state $(s \setminus \mathrm{del}(a)) \cup \mathrm{add}(a)$.

\textbf{\ac{pddl}} \cite{aeronautiques1998pddl} gives this a standardized, more expressive syntax, splitting a problem $P = \langle D, I \rangle$ into a domain $D$ and an instance $I$. The \textbf{domain} $D$ declares \textit{predicates} (relations such as \texttt{(on ?x ?y)}) and \textit{action schemas}---parameterized templates that generalize STRIPS operators, e.g.:
            \begin{lstlisting}[style=PDDL,  label=lst:pddl_stack_action]
                (:action stack
                    :parameters (?obj1 - block ?obj2 - block)
                    :precondition (and (holding ?obj1) (clear ?obj2))
                    :effect (and (on ?obj1 ?obj2) (not (holding ?obj1))
                                 (not (clear ?obj2)) (handempty) (clear ?obj1))
                )

            \end{lstlisting}
whose negative effects (\texttt{(holding ?obj1)}, \texttt{(clear ?obj2)}) form the delete-list and whose positive effects form the add-list. The \textbf{instance} $I$ lists the \textit{objects}, the \textit{initial state} $s_0$ (the ground atoms true at the start), and the \textit{goal} $g$ (ground atoms required in any goal state). Grounding the schemas---replacing their parameters by objects---yields the ground actions. PDDL also supports features beyond STRIPS, such as conditional effects (\texttt{(when <condition> <effect>)}), numeric fluents, and axioms (derived predicates).

This syntax induces a \textit{state model} $S(P) = \langle S, s_0, S_G, Act, A, f \rangle$ as its semantics \cite{geffner2013concise, hofmann2024learning}: $S$ is the set of all states; $s_0 \in S$ the initial state; $Act$ the set of ground actions; $A(s) = \{\, a \in Act : \mathrm{pre}(a) \subseteq s \,\}$ the actions applicable in $s$; $f(a, s) = (s \setminus \mathrm{del}(a)) \cup \mathrm{add}(a)$ the successor state for $a \in A(s)$; and $S_G = \{\, s \in S : g \subseteq s \,\}$ the goal states. A \textit{solution} (or \textit{plan}) is a sequence of actions $\sigma = [a_0, \dots, a_n]$ with $a_i \in A(s_i)$ and $s_{i+1} = f(a_i, s_i)$ that ends in a goal state $s_{n+1} \in S_G$; an \textit{optimal} plan minimizes length or an associated cost.
%Derived predicates allow the definition of new predicates whose truth values are logically inferred from the base predicates in the current state, rather than being directly changed by action effects. For example, one might define:
% \begin{lstlisting}[style=PDDL, language=Lisp, basicstyle=\ttfamily\scriptsize]
% (:derived (block_in_region ?b - block ?r - region)
%     (exists (?p - pose)
%         (and (at_pose ?b ?p)
%              (pose_in_region ?p ?r) ; Assuming pose_in_region is base or another derived
%         )
%     )
% )
% \end{lstlisting}
% This allows for more expressive and modular domain definitions.
\subsection{Classical Planning Algorithms and Fast Downward}
\label{subsec:planning-algorithms}
Given a compact \ac{pddl} encoding, a planner must search the state space it
implies for a path from $s_0$ to a goal state. Because that space grows
exponentially with the number of atoms, modern classical planners treat the
task as \emph{heuristic state-space search}: they expand states with an
informed search algorithm---greedy best-first search, or A\textsuperscript{*}
\cite{hart1968formal} when an optimal plan is required---guided by a heuristic
that estimates the remaining cost from each state to the goal
\cite{geffner2013concise}. The performance of such a planner depends almost
entirely on the quality of this heuristic.

What makes the approach \emph{domain-independent} is that these heuristics are
derived automatically from the PDDL description rather than supplied by hand.
The most influential family relaxes the task by ignoring delete-lists; because
solving even this relaxation optimally is hard, planners use tractable
approximations such as the max and additive heuristics and the FF heuristic,
which extracts a plan for the relaxed task and counts its actions
\cite{hoffmann2001ff, geffner2013concise}. A complementary family reasons over
\emph{landmarks}---facts that every valid plan must eventually make
true---and the LAMA planner combines landmark and relaxation heuristics in a
cost-based anytime search \cite{richter2010lama}.

\textbf{Fast Downward} \cite{helmert2006fast} is the classical planner used as
the search backend throughout this work; PDDLStream invokes it to solve each
symbolic planning problem once its streams have supplied the continuous values
that problem requires. Rather than operate on STRIPS propositions directly,
Fast Downward first translates the grounded task into a \emph{multi-valued}
representation (SAS\textsuperscript{+}), in which each variable takes exactly
one value from a mutually exclusive set instead of many independent Booleans.
From this encoding it builds a \emph{causal graph} that captures how the
variables depend on one another, and derives structure-aware heuristics (such
as the causal-graph heuristic) that guide a best-first search. This pairing of
an automatic multi-valued translation with structure-exploiting heuristics is
what lets Fast Downward scale to the grounded planning tasks that arise in
this work.

\section{Planning under Partial Observability}\label{sec:pod}
Classical planning's assumption of full observability is often violated in robotics. When an agent has incomplete information about the true state of the world, it must reason about its \textit{belief state}.

\subsection{Belief States, K-Literals, and Know-If Fluents}
A belief state represents the agent's current set of possible world states consistent with its history of actions and observations. We refer to the setting this thesis targets as \emph{Partially Observable Deterministic (POD)} planning: the world itself behaves deterministically and the only uncertainty is the agent's incomplete knowledge of the initial situation, which sensing resolves. POD builds on \cite{bonet2014flexible, albore2009translation, geffner2013concise}. The belief state $b$ is explicitly a set of possible states. Actions transition this set, and observations prune it.
A common way to reason about knowledge in POD planning is through K-literals \cite{bonet2011planning, geffner2013concise}. K-literals make the robot's knowledge about a proposition $p$ explicit. For example:
\begin{itemize}
    \item $K(p)$ signifies "$p$ is known to be true."
    \item $K(\neg p)$ signifies "$p$ is known to be false."
\end{itemize}
If neither $K(p)$ nor $K(\neg p)$ holds, the truth of $p$ is unknown. Separately, $p$ is "possibly true" whenever $K(\neg p)$ does not hold---which also covers the case where $p$ is known to be true.

K-literals are the basis of a \emph{translation} that compiles a partially observable problem into a classical one \cite{geffner2013concise}. Each literal $L$ of the original problem is replaced by the two fluents $K(L)$ and $K(\neg L)$~\cite{palacios2006compiling}; in the translated problem these are made true in the initial state only where the agent actually knows the corresponding value, and a literal whose value is unknown contributes neither. The translated problem therefore carries no uncertainty in its initial state and can be solved by an ordinary classical planner, which now searches for a plan that reaches a desired \emph{state of knowledge}---for example $K(\text{holding}(\text{target\_obj}))$, i.e.\ the agent knows it holds the target---rather than a desired world state. This thesis instead carries belief with \acp{kif}, the know-whether construct of Brenner and Nebel~\cite{brenner2009continual}. A Know-If fluent is weaker than a K-literal: rather than asserting that a fact is known true ($K(L)$) or known false ($K(\neg L)$), it records only that the agent knows \emph{whether} the fact holds---that its value is settled---while the value itself is carried by an ordinary fluent. The two are therefore not interchangeable: $K(L)$ commits to $L$ being true, whereas a \ac{kif} states only that $L$'s value is known and leaves that value to the underlying fact. We adopt \acp{kif} because the domain already stores where each object is, so it suffices to add one Know-If fluent per fact---whether that location is known---rather than two separate knowledge literals $K(L)$ and $K(\neg L)$ (\cref{sec:kif_belief}).

\subsection{Probabilistic vs. Deterministic Partial Observability}
POD planning differs from planning with \acp{pomdp} \cite{kaelbling1998planning}---a distinction that matters because we later compare our approach to POMDP-based TAMP. In POMDPs, the world is modeled as a \ac{mdp} with probabilistic transitions and observations, and the belief state is a probability distribution over world states. While highly expressive, POMDPs are computationally very demanding. POD planning assumes outcomes and observations are deterministic but unknown. This makes it more tractable than POMDPs, and it fits robotic problems where the uncertainty comes from not knowing the world rather than from actions behaving randomly.


\section{Task and Motion Planning (TAMP)}
\Ac{tamp} connects high-level task goals with a robot's actual physical movements \cite{garrett2021integrated}. The core problem is ensuring that a plan is physically possible. High-level planners often ignore geometry, leading to plans that seem correct but cannot be executed. Yet, including all geometric details from the start makes the problem too complex to solve. Combining high-level logic with real-world geometry is the central difficulty of TAMP, and frameworks like PDDLStream exist to bridge it.


\subsection{PDDLStream for TAMP}
PDDLStream \cite{garrett2018pddlstream} is a TAMP framework built to bridge this symbolic-continuous gap. It extends classical \ac{pddl} by introducing the core concept of \textbf{streams}.

\paragraph{Streams: declarative procedures for continuous parameters.}
A stream in PDDLStream is a procedure that generates or tests values for continuous parameters required by symbolic actions. These parameters might include object poses, robot configurations, grasp transformations, or motion trajectories. Streams have both a procedural and a declarative component:
\begin{itemize}
    \item \textbf{Procedural Component (Conditional Generator):} This is the actual code (e.g., a Python function) that, given some input values (which can be discrete symbols or other continuous values), produces a (potentially infinite) sequence of output values. For example, an \ac{ik} stream, given an object, its target pose, and a grasp, would generate a sequence of robot configurations that achieve that end-effector pose. A collision-checking stream, given a trajectory, would test if it is collision-free.
    \item \textbf{Declarative Component (Stream Specification):} This part is described in a PDDL-like syntax and informs  the symbolic planner what the stream does, what inputs it requires, what outputs it produces, and, crucially, what logical conditions (fluents) are associated with these inputs and outputs. This specification includes:
        \begin{itemize}
            \item \texttt{:inputs}: The typed parameters the stream procedure takes.
            \item \texttt{:domain}: A set of PDDL facts (preconditions) that must be true for the stream to be applicable for a given set of inputs. For example, an IK stream might require that a valid pose and grasp for the object are already known.
            \item \texttt{:outputs}: The typed parameters that the stream procedure generates.
            \item \texttt{:certified}: A set of PDDL facts that are guaranteed to be true if the stream successfully produces a set of outputs for given inputs. These "certified fluents" are added to the planner's state and represent new knowledge derived from the continuous domain. For instance, if an IK stream successfully finds a configuration \texttt{?q} for picking object \texttt{?obj} at pose \texttt{?p} with grasp \texttt{?g}, it might certify the fluent \texttt{(KinSolution ?obj ?p ?g ?q)}.
        \end{itemize}
\end{itemize}
PDDLStream employs a \textbf{lazy evaluation} strategy. Streams are not called exhaustively beforehand; instead, the symbolic planner calls them "on-demand" when it encounters an action whose preconditions require a continuous parameter to be instantiated or a continuous condition to be tested.

\paragraph{Example PDDLStream workflow.}
Consider a symbolic action \texttt{(pick ?obj ?p ?g ?q ?t)} (pick object \texttt{?obj} from pose \texttt{?p} using grasp \texttt{?g} via configuration \texttt{?q} along trajectory \texttt{?t}). For this action to be applicable:
\begin{enumerate}
    \item A stream might first be called to sample a stable \texttt{?p} for \texttt{?obj} on a surface. If successful, it certifies \texttt{(AtPose ?obj ?p)}.
    \item Then, another stream samples a grasp \texttt{?g} for \texttt{?obj}. If successful, it certifies \texttt{(Grasp ?obj ?g)}.
    \item An IK stream then takes \texttt{?obj, ?p, ?g} as input and tries to find a robot configuration \texttt{?q}. If successful, it certifies \texttt{(KinSolution ?obj ?p ?g ?q)}.
    \item Finally, a motion planning stream takes the current configuration and \texttt{?q} as input to find a collision-free trajectory \texttt{?t}, certifying \texttt{(Motion ?q\_start ?t ?q)}.
\end{enumerate}
Only if all these streams succeed and certify their respective fluents will the preconditions of the symbolic \texttt{(pick ...)} action be met.

The overall PDDLStream algorithm often reduces a TAMP problem to a sequence of finite PDDL problems, iteratively calling streams to discover more certified fluents until a complete, grounded plan is found. The original PDDLStream framework, however, was primarily designed for fully observable and deterministic settings.

\section{Spatial Representation for Robotic Planning}
\label{sec:spatial_representation}

A robot needs some way to represent the 3D space around it and the objects in it. Representing space is harder under partial observability, because the robot must also represent what it does not yet know.

\subsection{Discretization of Continuous Space}
Robots operate in continuous space, but symbolic planners reason over discrete states --- so the continuous space has to be discretized somehow.

\paragraph{Voxel grids.}
A common approach to represent 3D environments is to discretize them into a grid of uniform volumetric pixels, or \textbf{voxels}. Each voxel can store information, such as whether it is occupied, free, or unknown. While straightforward, uniform voxel grids can be memory-intensive and computationally expensive for large environments or high resolutions.

\textbf{Octrees} (\cref{fig:octmap_illustration}) \cite{hornung2013octomap} offer a more efficient hierarchical data structure for representing 3D space. An octree is a tree data structure in which each internal (non-leaf) node has exactly eight children, while leaf nodes have none. It starts with a single large voxel encompassing the entire space of interest. If this voxel is not homogeneous (e.g., it contains both free and occupied space), it is recursively subdivided into eight smaller child voxels (octants).
\begin{figure}[H]
    \centering % Centers the image
    \includegraphics[width=0.65\textwidth]{octmap_illustration.png}
    \caption{Example of an octree storing free (shaded white) and occupied (black) cells. The volumetric model is shown on the left and the corresponding tree representation on the right.}
    \label{fig:octmap_illustration}
\end{figure}
This subdivision continues until voxels are homogeneous or a minimum resolution/maximum depth is reached. This process creates a tree where large, uniform regions (whether free, occupied, or unknown) are represented by single nodes at a shallower depth, making it highly efficient for representing sparse environments in comparison to uniform grid voxelization. The \textbf{OctoMap} framework \cite{hornung2013octomap} uses this principle to build probabilistic 3D occupancy maps from noisy sensor data. In such a framework, unobservable space is typically represented implicitly by the absence of nodes in a given volume; querying such a region returns an "unknown" status. The same octree structure also underpins learning-based 3D methods such as OctNet \cite{riegler2017octnet}, which runs deep convolutions efficiently on an octree rather than building an occupancy map; here it is the occupancy-mapping use that is relevant.

% Octrees are widely used in robotics for creating 3D occupancy maps from sensor data (e.g., OctoMap \cite{hornung2013octomap}), as they efficiently represent large free or unknown areas with fewer nodes.

% However, a limitation of standard Octree construction methods, when applied to planning under partial observability, is their typical reliance on direct sensor measurements (like point clouds) to guide subdivision based on observed point density or occupancy changes. This makes it challenging to use them directly to represent and reason about beliefs in unobserved or occluded regions where specific target objects might be hiding – a key focus for information-gathering plans.

\paragraph{Critical Regions.}
To make planning more efficient, space can be divided based on task-relevant meaning rather than uniform geometric criteria. One such approach is the concept of \textbf{\acp{cr}}, introduced by Molina et al.~\cite{molina2020learn} and later extended by Shah and Srivastava~\cite{shah2022abstractions}. In the original formulation, a critical region is a part of the configuration space with a low probability of being sampled by a uniform sampler yet essential to a solution---narrow passages such as doorways~\cite{molina2020learn}. Shah and Srivastava generalised this to regions where task-relevant features---controller success rates, object visibility, or transition likelihoods---change significantly~\cite{shah2022abstractions}.

In this line of work the regions are not hand-specified: a deep network is trained to \emph{predict} the critical regions of a given environment from that scene's input, and the predictor generalises to environments it was not trained on. Once predicted, a \ac{rbvd} partitions the configuration space around the CRs: each point is assigned to the nearest critical region, forming abstract states. Transitions between these abstract states define abstract actions, forming the backbone for hierarchical planning. By focusing computation on these task-relevant regions, this approach scales planning more effectively. In this formulation the regions are predicted once per environment and then held fixed, so they do not adapt as objects move within a scene during execution; later work in the same line learns abstractions that adapt as the task changes~\cite{nayyar2025option}, though for reinforcement learning rather than POD-TAMP.

Our work borrows the idea of task-relevant spatial regions, but constructs Boxels geometrically from detected objects and their occlusions rather than learning a region predictor (\cref{ch:methods}).
